\documentclass[12pt]{article}

\usepackage[utf8]{inputenc}
\usepackage[margin=1in]{geometry}
\usepackage{setspace}
\usepackage{url}
\usepackage{mathptmx}

\doublespacing
\setlength{\emergencystretch}{3em}

\title{Reading Canada Through Short Fiction}

\author{%
  Ma Toan Bach\\
  Student Number: 112708227\\
  LSO 100--Canadian Short Story (KCC)\\
  Final Assignment
}

\date{August 11, 2026}

\begin{document}

\maketitle

\section*{Part 1 -- Looking Back: Your Journey Through Canadian Short Fiction}

At the beginning of the course, I assumed a Canadian short story would be defined mainly by setting: small towns, winter, wilderness, or recognizable national details. After reading across the semester, I understand Canadian short fiction in a more complicated way. The stories are not only about where Canada is located; they are about who is allowed to belong, whose grief is heard, and how place can be beautiful and damaging at the same time.

Margaret Laurence's ``The Loons'' helped me see this first. The story connects landscape to colonial loss through Piquette and through the disappearance of the loons. The lake setting is not just scenery; it reminds readers that Indigenous presence can be pushed aside while settlers imagine the place as peaceful or natural (Laurence, 1970). This forced me to reconsider landscape as a record of exclusion, memory, and unequal power.

Bharati Mukherjee's ``The Management of Grief'' expanded that idea in a different direction. The story begins with a public Canadian tragedy, but its real focus is on how immigrant families grieve inside and against official systems. Shaila's grief is shaped by family loss, cultural expectations, and the pressure to become useful to others before she has fully understood her own pain (Mukherjee, 1988). This story showed me that multicultural Canada is not only a celebration of difference; it also includes misunderstanding, loneliness, and emotional labour.

Bridget Canning's ``Pack Ice Season'' brought the course into contemporary Canada. The pack ice becomes an event that people watch, discuss, and interpret, so the environment is not background but a force that changes how a community is seen by others (Canning, 2020). That detail made me reconsider the idea that contemporary Canadian stories move away from landscape; instead, they show landscape being reshaped by media, tourism, climate anxiety, and public attention. Together, these stories reveal why Canada cannot be represented by one national experience. Indigenous loss, immigrant mourning, and environmental uncertainty are different histories, but each challenges a simple image of Canada as peaceful or automatically inclusive.

\section*{Part 2 -- Making Connections: What Do These Stories Say About Canada?}

Alice Munro's ``Boys and Girls,'' Thomas King's ``Borders,'' and Roch Carrier's ``The Hockey Sweater'' all show that belonging in Canada is never neutral. Each story focuses on a different pressure: gender expectation, colonial state categories, and cultural-national loyalty. Together, they suggest that recognition often depends on rules created by families, governments, communities, or national symbols. Through point of view, symbolism, and irony, the three stories reveal identity as a conflict between self-understanding and labels imposed from outside.

All three stories present belonging as conditional, but the conditions come from different sources. In ``Boys and Girls,'' the narrator first feels recognized when her father introduces her as his ``new hired hand,'' making her feel ``red in the face with pleasure'' (Munro, 1968, p.~5). Yet the final judgment, ``She's only a girl,'' reduces her action and identity to gender (Munro, 1968, p.~16). Munro shows a child gradually learning a role prepared for her. King's ``Borders'' moves the conflict from family to state. The mother is asked to identify herself as Canadian or American, but repeatedly answers Blackfoot, even when that refusal leaves the family stuck between borders (King, 1993). Unlike Munro's narrator, she does not accept the imposed category. Carrier's ``The Hockey Sweater'' shows a third version: the narrator is judged by community symbols when he is forced to wear the wrong team's sweater (Carrier, 1984). Belonging can be denied quietly through gender, officially through paperwork, or socially through public shame.

The authors also use point of view to show how identity is learned, resisted, or misunderstood. Munro's retrospective narration lets the reader feel childhood pride while seeing what the older narrator understands: the father's praise was limited by gender. King also uses a child narrator, but the child does not fully control the political meaning of the mother's refusal, so the reader recognizes the dignity of her answer before the narrator explains it. Carrier's child-centered narration turns cultural pressure into comedy, but the humour depends on humiliation. When the narrator wears the Toronto Maple Leafs sweater in a community devoted to the Montreal Canadiens, childish disappointment becomes a lesson in French-Canadian communal identity (Carrier, 1984). In all three stories, childhood perspective makes belonging immediate rather than abstract.

Symbolism deepens this comparison because each story turns an ordinary object or place into a test of recognition. Munro's fox farm represents usefulness, freedom, and closeness to the father, but it also reveals that the narrator's access to masculine work is temporary. King's border makes colonial national categories visible and shows how they fail to recognize Indigenous sovereignty. Carrier's sweater is comic, but it makes communal loyalty visible on the narrator's body. A farm, a border booth, and a hockey sweater are concrete, but each carries a larger question: who is allowed to name themselves, and who must accept the role assigned by others?

The stories differ most sharply in how characters respond to pressure. Munro's narrator resists briefly when she opens the gate for Flora, but the ending suggests that she begins to internalize the word ``girl'' as a limit. King's mother refuses the official choice altogether, grounding identity in Blackfoot nationhood rather than Canadian or American paperwork. Carrier's narrator makes no political argument, but his embarrassment shows how strongly children absorb regional and linguistic loyalty before they can explain it. Gendered family roles, Indigenous sovereignty, and French-Canadian cultural loyalty all shape belonging differently.

Considered together, these stories make Canada appear less like a finished identity and more like a set of contested relationships. Munro shows belonging limited by patriarchy inside the home. King shows Indigenous identity exceeding the settler state's map. Carrier shows how language, region, religion, and sport can unite a community while punishing difference. The stories complicate the idea that Canada is simply inclusive: inclusion always has terms, and short fiction can reveal those terms through small, memorable moments.

\section*{Part 3 -- Looking Forward: The Story That Should Stay}

If a future LSO 100 class could read only one story from the course, I believe Madeleine Thien's ``Simple Recipes'' should remain on the syllabus. The story deserves to stay because it brings together many central questions in a compact and emotionally powerful way: immigration, family, cultural inheritance, memory, violence, and belonging. It teaches something especially valuable: Canadian literature can examine national questions through the private space of a kitchen, where love, authority, language, food, and shame are connected.

One important issue the story explores is multicultural life inside the home. Public discussions of multiculturalism often focus on festivals, food, language, and tolerance, but Thien makes the issue more intimate. The father's cooking initially represents care and continuity. The narrator remembers learning the ``simple recipe for making rice'' and watching her father remove ``tiny imperfections'' from the grains (Thien, 2002, p.~1). This image presents inherited culture as something taught through touch, repetition, and attention. However, the dinner-table violence later changes the meaning of food. A meal that should bring the family together becomes connected to shame, anger, and fear. Future students should read this story because it shows how inheritance can be loving and painful at the same time.

The first literary technique future students should study is symbolism. Rice, fish, and recipes carry more than one meaning. Rice symbolizes care, order, and inheritance when the father teaches the narrator. Fish symbolizes tension when the brother refuses the meal and the father reads that refusal as betrayal. The title itself, ``Simple Recipes,'' becomes ironic because family life is not simple at all. A recipe suggests that correct steps produce a reliable result. Thien's story shows the opposite: immigration, family loyalty, and cultural inheritance do not come with stable instructions. The title gives students a clear symbol for the impossibility of creating a simple formula for belonging.

The second important technique is retrospective first-person narration. The narrator tells the story from memory, so she can hold contradictory feelings at the same time. She remembers her father as gentle and careful, but she also remembers his violence. Near the end, she says that ``this violence will turn all my love to shame and grief'' (Thien, 2002, p.~5). The sentence refuses an easy judgment. The narrator does not simply stop loving her father, but she cannot love him innocently anymore. The retrospective voice presents trauma as a memory that continues to reshape love afterward.

Future students should pay particular attention to the dinner-table scene. The most important issue is cultural loyalty enforced through paternal authority. The brother's refusal of the fish becomes more than a complaint about food because the father treats it as rejection of the family and of what the family has carried into Canada. Masculinity, gratitude, and assimilation contribute to that conflict: the father expects obedience, demands thankfulness, and reacts as if refusal means turning away from one's origins. The scene transforms the kitchen from care into fear. It also forces the narrator into a painful position: she loves her father, but witnessing his violence makes that love morally dangerous. Few course stories show so clearly how immigration can become a crisis inside one family meal.

``Simple Recipes'' also connects meaningfully to Munro's ``Boys and Girls.'' Both stories show fathers shaping identity through work, skill, and authority, but the authority works differently. Munro's narrator feels proud when her father calls her his ``new hired hand,'' then later has that belonging reduced by ``only a girl'' (Munro, 1968, pp.~5, 16). Her father's power is quiet and social: it confirms the gender rule she has been learning. Thien's father also teaches through skill, but his authority becomes openly violent, dividing the narrator's memory between tenderness and fear. The comparison helps students see that family can enforce identity through ordinary language and physical force.

The story also connects to Mukherjee's ``The Management of Grief.'' Shaila is treated as useful because she seems calm after the plane disaster, even though her private grief remains unresolved (Mukherjee, 1988). Thien's narrator faces a different situation, but both stories show pain that cannot be easily translated into public language. In Mukherjee, grief is shaped by institutions and community expectations after catastrophe. In Thien, grief is shaped by a domestic scene the narrator cannot separate from love. Together, they show immigrant experience beyond arrival or success: silence, divided loyalty, and emotional pressure.

After completing this course, ``Simple Recipes'' helped me understand Canadian literature differently because it shows Canada through the emotional details of one family. It proves that Canadian experience can be found in a bowl of rice, a dinner table, a remembered sentence, or a silence after violence. That is why the story should remain on the syllabus: it teaches students to read Canadian short fiction as personal, cultural, and historical all at once.

\clearpage

\begin{center}
\textbf{References}
\end{center}

\hangindent=0.5in\noindent Canning, B. (2020). Pack ice season. In \textit{Some people's children}. Breakwater Books.

\hangindent=0.5in\noindent Carrier, R. (1984). \textit{The hockey sweater} (S. Fischman, Trans.). Tundra Books.

\hangindent=0.5in\noindent King, T. (1993). Borders. In \textit{One good story, that one}. HarperCollins.

\hangindent=0.5in\noindent Laurence, M. (1970). The loons. In \textit{A bird in the house}. McClelland \& Stewart.

\hangindent=0.5in\noindent Mukherjee, B. (1988). The management of grief. In \textit{The middleman and other stories}. Grove Press.

\hangindent=0.5in\noindent Munro, A. (1968). \textit{Boys and girls} [PDF]. \url{https://jerrywbrown.com/wp-content/uploads/2014/02/Boys-and-Girls-by-Alice-Munro.pdf}

\hangindent=0.5in\noindent Thien, M. (2002). \textit{Simple recipes} [PDF]. MostlyFiction. \url{http://mostlyfiction.com/excerpts/simplerecipes.htm}

\end{document}
